\section{Discussion}
\label{sec:discussion}

\paragraph{Why Qwen's best gate is also the learner-NWP gate.}
\citet{qiu2025gated} attribute the dominance of $G_1$ over $G_2$--$G_5$
to (i) non-linearity between the value and output projections and (ii)
input-dependent sparsity that mitigates the attention sink. Both
properties have a clean interpretation in the L2-NWP setting: a
proficiency-aware gate needs to selectively suppress or amplify the
pieces of context a learner \emph{actually uses} when predicting the
next word. An A1 Spanish-L1 learner writing in English uses the
immediate left context heavily and the distal discourse context
barely; a C2 learner integrates more globally. The gate at $G_1$ is
the place in the attention block where that filtering has the most
direct effect, because it operates on the softmax-weighted value
aggregate --- the attention output --- and not on the pre-aggregation
$Q/K/V$ projections. In other words, we expect the ordering
$G_1 > G_2 > \{G_3, G_4\} > G_5$ to reproduce on learner NWP, with a
wider margin than on native LM benchmarks because the operation is
more load-bearing here.

\paragraph{Relation to the CMCL companion finding.}
The companion paper showed that a data-only learner GPT-2 produces a
\emph{selectively-shifted} error profile --- not uniformly more errors,
but systematically more errors in the SLA-predicted vulnerable
categories \citep{stearns2026artificial}. Our proposal is the
architectural next step: if data-only continual pre-training bends the
error profile in the right direction for the EFCAMDAT average, an
architectural gate with metadata input should bend it per learner
profile. The CMCL result is the existence proof for the data-only
version of this claim; this paper is the control knob. Quantitative
validation via the ERRANT-based error-histogram pipeline is a
well-scoped follow-up (see Future Work); in the present pilot the
two metrics --- stratified PPL and cloze accuracy --- already
constrain the hypothesis meaningfully.

\paragraph{The perplexity-as-diagnostic connection.}
\citet{stearns2025nwpaleval} argues for treating perplexity not as an
end-to-end metric but as a diagnostic that identifies linguistic
constructions a learner finds uncertain. Our gate is the
minimum-intervention operationalisation: the gating scores themselves
are a layer-by-layer, head-by-head record of which contextual
information the model allows to flow through when conditioned on a
given learner profile. Inspecting those scores on a fixed input while
varying the metadata (A1 Spanish vs C1 Japanese) provides a principled
way of extracting \emph{what distinguishes these learners' expectations}
directly from the model's internal mechanisms, rather than from
downstream behavioural probes.

\paragraph{What the variant catalogue is for.}
The nine variants in \texttt{experiments-ideas/} are not ablations
written post-hoc to defend the main result. They are pre-registered
tests of the specific causal story we are making: that gate site
matters in the order Qwen predicts (v02), that elementwise granularity
matters but perhaps not as much as assumed (v03), that the
multiplicative form beats the additive (v04), that both metadata axes
contribute (v05, v06), that typology generalises where L1 ID does not
(v07), that the same mechanism helps cloze (v08), and that the
backbone does not need to move for the gate to do its job (v09). Each
variant has a stand-alone \texttt{.md} that specifies its expected
outcome \emph{before the experiment is run}; the flagship run either
produces results consistent with that causal map or forces us to
revise it.
